
\documentclass{article}
\usepackage{geometry}
\geometry{margin=1.5cm, vmargin={0pt,1cm}}
\setlength{\topmargin}{-1cm}
\setlength{\headheight}{12.5pt}
\setlength{\paperheight}{29.7cm}
\setlength{\textheight}{25.3cm}

% useful packages.
% \usepackage[UTF8]{ctex}
\usepackage{fancyhdr}
\usepackage{layout}
\usepackage{tikz}
\usetikzlibrary{arrows.meta}

% import common macros
% % % % % % % % % % % % % % % % % % % % % % % % % % % % % % % %
% Common Macros for LaTeX
% % % % % % % % % % % % % % % % % % % % % % % % % % % % % % % %

% Useful packages
\usepackage{amsfonts}
\usepackage{amsmath}
\usepackage{amssymb}
\usepackage{amsthm}
\usepackage{enumerate}
\usepackage{graphicx}
\usepackage{multicol}
\usepackage[ruled,linesnumbered]{algorithm2e}

% Math operators and common symbols
\newcommand{\dif}{\mathrm{d}}
\newcommand{\innerProd}[1]{\left\langle #1 \right\rangle}
\newcommand{\difFrac}[2]{\frac{\dif #1}{\dif #2}}
\newcommand{\pdfFrac}[2]{\frac{\partial #1}{\partial #2}}
\newcommand{\T}{\mathrm{T}}
\newcommand{\norm}[1]{\left\| #1 \right\|}
\newcommand{\abs}[1]{\left| #1 \right|}

% Vector and Matrix shortcuts
\newcommand{\vecTwo}[2]{
  \begin{bmatrix}
    #1 \\
    #2
\end{bmatrix}}
\newcommand{\vecThree}[3]{
  \begin{bmatrix}
    #1 \\
    #2 \\
    #3
\end{bmatrix}}
\newcommand{\vecTwoH}[2]{
  \begin{bmatrix}
    #1 & #2
  \end{bmatrix}
}
\newcommand{\vecThreeH}[3]{
  \begin{bmatrix}
    #1 & #2 & #3
  \end{bmatrix}
}

% Common fields
\newcommand{\R}{\mathbb{R}}
\newcommand{\C}{\mathbb{C}}
\newcommand{\Q}{\mathbb{Q}}
\newcommand{\Z}{\mathbb{Z}}
\newcommand{\N}{\mathbb{N}}

% Common Operators
\renewcommand{\Re}{\mathrm{Re}}
\renewcommand{\Im}{\mathrm{Im}}

% Theorem-like environments
\theoremstyle{definition}
\newtheorem{theorem}{Theorem}[section]
\newtheorem{lemma}[theorem]{Lemma}
\newtheorem{corollary}[theorem]{Corollary}
\newtheorem{definition}[theorem]{Definition}
\newtheorem{remark}[theorem]{Remark}
\newtheorem{exercise}{Exercise}[section]

% Support for individual Chinese characters using fontspec (for LuaLaTeX)
\usepackage{fontspec}
\newfontfamily\zhfont{SimSun} % Search system fonts by name
\newcommand{\zhstr}[1]{{\zhfont #1}}

% Solutions and proofs
\AtBeginDocument{
  \ifdefined\ctexset
  \newenvironment{solution}{
  \begin{proof}[解]}{
  \end{proof}}
  \else
  \newenvironment{solution}{
  \begin{proof}[Solution]}{
  \end{proof}}
  \fi
}

\usepackage{luacode}
% Utility to get environment variables (requires lualatex)
\newcommand{\getEnv}[2]{\directlua{tex.print(os.getenv("#1") or "#2")}}


% Name and uID
\newcommand{\name}{\getEnv{NAME}{NAME}}
\newcommand{\uid}{\getEnv{UID}{UID}}
\newcommand{\course}{Complex Analysis}
\newcommand{\hwNum}{3}

\begin{document}

\pagestyle{fancy}
\fancyhead{}
\lhead{\zhstr{\name} (\uid)}
\chead{\course\ \#\hwNum}
\rhead{\today}

\section{Exercise C.335}

\begin{theorem}[Theorem C.334]
  The exponential of a skew-Hermitian matrix is a unitary matrix.
  Conversely, any unitary matrix must be the exponential of some
  skew-Hermitian matrix.
\end{theorem}

\begin{proof}
  Consider a skew-Hermitian matrix $A$ with $A^*=-A$. Using the power
  series, we have
  \[
    (e^A)^* = \left( \sum_{k=0}^\infty \frac{1}{k!} A^k\right)^*
    = \sum_{k=0}^\infty \frac{1}{k!} (A^*)^k
    = \sum_{k=0}^\infty \frac{1}{k!} (-A)^k
    = e^{-A}.
  \]
  Then we need to show that $A$ and $-A$ are commutative, which is
  true since $A$ and $A^*$ are always commutative. Hence
  \[
    (e^A)^*(e^A) = e^{-A} e^A = I.
  \]
  Therefore, $e^A$ is unitary.

  From another perspective, let $U$ be a unitary matrix. Then
  $U^*U=I$, so $U^*$ is the inverse of $U$. We need to show that there exists
  a skew-Hermitian matrix $A$ such that $e^A=U$.

  Consider unitary diagonalization of $U$:
  \[
    U = V \Lambda V^*,
  \]
  where $V$ is unitary and $\Lambda$ is diagonal with entries on the
  unit circle. Write $\Lambda = \mathrm{diag}(e^{i\theta_1}, \ldots,
  e^{i\theta_n})$ with $\theta_j \in \R$. Then we define $A$ as:
  \[
    A = V \mathrm{diag}(i\theta_1, \ldots, i\theta_n) V^*.
  \]
  Since $\mathrm{diag}(i\theta_1, \ldots, i\theta_n)$ is
  skew-Hermitian, $A$ is also skew-Hermitian. Moreover, using the power
  series expansion of the exponential function, we have
  \[
    e^A = V e^{\mathrm{diag}(i\theta_1, \ldots, i\theta_n)} V^* = V
    \Lambda V^* = U.
  \]
  This completes the proof.

\end{proof}

\section{Exercise 2.12}

The function
\[
  z \mapsto \Arg z
\]
is continuous on the complex plane slit along the negative real axis
\[
  \C^- := \C \setminus \{t \in \R : t \le 0\}.
\]

\begin{proof}
  Let $z_0 = x_0 + iy_0 \in \C^-$ with $z_0 \ne 0$.
  We need to show $\Arg z \to \Arg z_0$ as $z \to z_0$.

  Write $z = x + iy = re^{i\theta}$ where $r = |z|$, $\theta = \Arg z
  \in (-\pi, \pi)$,
  and similarly $z_0 = r_0 e^{i\theta_0}$.

  \textbf{Case 1: $x_0 > 0$ (right half-plane).}

  Here $\Arg z_0 = \arctan(y_0/x_0)$.
  For $z$ sufficiently close to $z_0$, we have $x > 0$, so
  \[
    \Arg z = \arctan\frac{y}{x}.
  \]
  Since $(x,y) \mapsto y/x$ is continuous at $(x_0,y_0)$ with $x_0 > 0$,
  and $\arctan$ is continuous on $\R$, the composition is continuous.
  Hence $\Arg z \to \Arg z_0$ as $z \to z_0$.

  \textbf{Case 2: $x_0 < 0$ and $y_0 > 0$ (upper left quadrant).}

  Here $\Arg z_0 = \arctan(y_0/x_0) + \pi \in (0,\pi)$.
  For $z$ close to $z_0$, we have $x < 0$ and $y > 0$, so
  \[
    \Arg z = \arctan\frac{y}{x} + \pi.
  \]
  Again this is a composition of continuous functions, so $\Arg$ is continuous.

  \textbf{Case 3: $x_0 < 0$ and $y_0 < 0$ (lower left quadrant).}

  Here $\Arg z_0 = \arctan(y_0/x_0) - \pi \in (-\pi, 0)$.
  For $z$ close to $z_0$, we have $x < 0$ and $y < 0$, so
  \[
    \Arg z = \arctan\frac{y}{x} - \pi,
  \]
  which is continuous.

  \textbf{Case 4: $x_0 = 0$ and $y_0 > 0$ (positive imaginary axis).}

  Here $\Arg z_0 = \pi/2$. Given $\varepsilon > 0$, we need $\delta >
  0$ such that
  $|z - z_0| < \delta$ implies $|\Arg z - \pi/2| < \varepsilon$.

  Since $\cot\theta = x/y$ for $\theta \in (0,\pi)$, we have
  \[
    |\Arg z - \pi/2| = |\arccot(x/y)|.
  \]
  Using the inequality $|\arccot t| \le |t|$ for $t$ near $0$ (since
  $\arccot'(0) = -1$),
  for $|x/y| < \min(\varepsilon, 1)$ we have
  \[
    |\Arg z - \pi/2| \le \left|\frac{x}{y}\right|.
  \]
  Choose $\delta = y_0 \cdot \min(\varepsilon, 1)/2$. Then $|z - z_0|
  < \delta$ implies
  $|x| < \delta$ and $y > y_0 - \delta \ge y_0/2$, so
  \[
    \left|\frac{x}{y}\right| < \frac{\delta}{y_0/2}
    = \frac{2\delta}{y_0}
    \le \min(\varepsilon, 1) \le \varepsilon.
  \]

  \textbf{Case 5: $x_0 = 0$ and $y_0 < 0$ (negative imaginary axis).}

  Here $\Arg z_0 = -\pi/2$. The argument is symmetric to Case 4, using
  $|\Arg z + \pi/2| = |\arccot(x/y)|$ with $\cot(-\pi/2 + \theta) =
  -\tan\theta$.

  \textbf{Case 6: $z_0$ on the positive real axis ($y_0 = 0$, $x_0 > 0$).}

  Here $\Arg z_0 = 0$. For $z = x + iy$ close to $z_0$,
  \[
    |\Arg z| = |\arctan(y/x)|.
  \]
  Using $|\arctan t| \le |t|$ for $t \in \R$,
  \[
    |\Arg z| \le \left|\frac{y}{x}\right|.
  \]
  For $|z - z_0| < \delta$ with $\delta < x_0/2$, we have $x > x_0/2$, so
  \[
    |\Arg z| \le \frac{|y|}{x_0/2} < \frac{2\delta}{x_0}.
  \]
  Choosing $\delta = \min(x_0/2, x_0\varepsilon/2)$ gives $|\Arg z| <
  \varepsilon$.

  In all cases, $\Arg z$ is continuous at each $z_0 \in \C^-$.
\end{proof}

\section{Exercise 2.21}

\begin{theorem}[Theorem 2.20]
  Let $f,g : D \to \C$ be complex differentiable at $a \in D$ and
  let $\lambda \in \C$ with $f(a) \ne 0$. Then
  $\lambda f$, $f+g$, $fg$, and $1/f$ are complex differentiable at $a$, and
  \[
    (\lambda f)'(a) = \lambda f'(a),
    \qquad
    (f+g)'(a) = f'(a) + g'(a),
  \]
  \[
    (fg)'(a) = f'(a)g(a) + f(a)g'(a),
    \qquad
    \left(\frac{1}{f}\right)'(a) = -\frac{f'(a)}{f(a)^2}.
  \]
\end{theorem}

Prove Theorem 2.20.

\begin{proof}
  Let $a\in D$.

  For scalar multiplication,
  \[
    \frac{(\lambda f)(z)-(\lambda f)(a)}{z-a}
    = \lambda\frac{f(z)-f(a)}{z-a}
    \xrightarrow[z\to a]{} \lambda f'(a).
  \]
  So $(\lambda f)'(a)=\lambda f'(a)$.

  For addition,
  \[
    \frac{(f+g)(z)-(f+g)(a)}{z-a}
    = \frac{f(z)-f(a)}{z-a} + \frac{g(z)-g(a)}{z-a}
    \xrightarrow[z\to a]{} f'(a)+g'(a).
  \]
  So $(f+g)'(a)=f'(a)+g'(a)$.

  For product,
  \[
    \begin{aligned}
      \frac{f(z)g(z)-f(a)g(a)}{z-a}
      &= \frac{f(z)g(z)-f(z)g(a)}{z-a}
      + \frac{f(z)g(a)-f(a)g(a)}{z-a} \\
      &= f(z)\frac{g(z)-g(a)}{z-a}
      + g(a)\frac{f(z)-f(a)}{z-a}.
    \end{aligned}
  \]
  Because differentiability implies continuity, $f(z)\to f(a)$ as
  $z\to a$. Taking limits gives
  \[
    (fg)'(a)=f(a)g'(a)+g(a)f'(a).
  \]

  For reciprocal, assume $f(a)\neq 0$. By continuity of $f$, $f(z)\neq 0$
  for $z$ near $a$, and
  \[
    \frac{\frac1{f(z)}-\frac1{f(a)}}{z-a}
    = \frac{f(a)-f(z)}{(z-a)f(z)f(a)}
    = -\frac{1}{f(z)f(a)}\frac{f(z)-f(a)}{z-a}.
  \]
  Taking limits,
  \[
    \left(\frac1f\right)'(a)
    = -\frac{1}{f(a)^2}f'(a)
    = -\frac{f'(a)}{f(a)^2}.
  \]
\end{proof}

\section{Exercise 2.22}

Show that the function $f : D \to \C$ given by
\[
  f(z) = z^n,
\]
where
\[
  D = \C \text{ if } n \in \N,
  \qquad
  D = \C^\bullet \text{ if } n \in -\N_+,
\]
is complex differentiable and
\[
  f'(z)=n z^{n-1}.
\]
Then deduce that the complex derivative of the polynomial
\[
  p_n(z)=\sum_{j=0}^{n} a_j z^j
\]
is
\[
  p_n'(z)=\sum_{j=1}^{n} j a_j z^{j-1}.
\]

\begin{proof}
  First let $n\in\N$. For $n=0$, $f(z)=1$ and $f'(z)=0$, which matches
  $nz^{n-1}=0$. For $n=1$, $f(z)=z$ and $f'(z)=1$.

  Assume $(z^k)'=kz^{k-1}$. Then by the product rule,
  \[
    (z^{k+1})'=(z\cdot z^k)'=z'\,z^k+z\,(z^k)'=z^k+kz^k=(k+1)z^k.
  \]
  So by induction, $(z^n)'=nz^{n-1}$ for all $n\in\N$.

  Now let $n=-m$ with $m\in\N_+$. Then on $\C^\bullet$,
  \[
    z^n=z^{-m}=\frac{1}{z^m}.
  \]
  Using the reciprocal rule (Theorem 2.20) and the positive-power result,
  \[
    (z^{-m})'=-\frac{(z^m)'}{(z^m)^2}
    =-\frac{mz^{m-1}}{z^{2m}}
    =-mz^{-m-1}
    =nz^{n-1}.
  \]
  Hence for every integer exponent in the stated domain,
  \[
    (z^n)'=nz^{n-1}.
  \]

  Finally, for a polynomial $p_n(z)=\sum_{j=0}^n a_j z^j$, linearity of
  derivative gives
  \[
    p_n'(z)=\sum_{j=0}^n a_j (z^j)'
    =\sum_{j=1}^n j a_j z^{j-1}.
  \]
\end{proof}

\section{Exercise 2.42}

Give a visual differentiation of the principal branch of the logarithm
function defined in Theorem 1.109.
You must draw two plots similar to those in Example 2.41 and elaborate on
how the derivative is derived from the plots.

\begin{solution}
  Write
  \[
    z_0=r_0e^{i\theta_0},\qquad w_0=\Log z_0=\ln r_0+i\theta_0,
    \qquad r_0>0,\ \theta_0\in(-\pi,\pi).
  \]
  We draw two local experiments, in the style of Example 2.41.

  \textbf{Plot A (angular increment, radius fixed).}

  \begin{center}
    \begin{minipage}{0.44\textwidth}
      \centering
      \begin{tikzpicture}[scale=1.0,>=Stealth]
        \draw[->] (-0.2,0) -- (3.0,0) node[right] {$\Re z$};
        \draw[->] (0,-0.2) -- (0,2.6) node[above] {$\Im z$};
        \draw[thick] (-2.2,0) -- (-0.15,0);
        \node[above] at (-1.3,0) {slit};
        \draw (0,0) -- (1.9,1.1) node[midway,below right] {$r_0$};
        \draw (0,0) -- (1.55,1.55);
        \draw[dashed] (1.9,1.1) arc (30:45:2.2);
        \fill (1.9,1.1) circle (1pt) node[right] {$z_0$};
        \fill (1.55,1.55) circle (1pt);
        \draw[->] (1.87,1.14) -- (1.60,1.50) node[above] {$\delta z$};
        \draw (0.45,0) arc (0:30:0.45);
        \node at (0.62,0.20) {$\theta_0$};
        \draw (0.62,0) arc (0:45:0.62);
        \node at (0.96,0.42) {$\theta_0+\delta\theta$};
      \end{tikzpicture}
    \end{minipage}
    \hfill
    \begin{minipage}{0.12\textwidth}
      \centering
      \begin{tikzpicture}[>=Stealth]
        \draw[->,thick] (0,0) -- (1.2,0) node[midway,above] {$\Log$};
      \end{tikzpicture}
    \end{minipage}
    \hfill
    \begin{minipage}{0.40\textwidth}
      \centering
      \begin{tikzpicture}[scale=1.0,>=Stealth]
        \draw[->] (-0.2,0) -- (3.0,0) node[right] {$u$};
        \draw[->] (0,-0.2) -- (0,2.6) node[above] {$v$};
        \draw[dashed] (1.6,0) -- (1.6,2.4);
        \fill (1.6,1.1) circle (1pt) node[right] {$w_0$};
        \fill (1.6,1.55) circle (1pt);
        \draw[->] (1.6,1.13) -- (1.6,1.50) node[right] {$i\,\delta\theta$};
        \node[below] at (1.6,0) {$\ln r_0$};
        \node[left] at (0,1.1) {$\theta_0$};
        \node[left] at (0,1.55) {$\theta_0+\delta\theta$};
      \end{tikzpicture}
    \end{minipage}
  \end{center}

  \textbf{Plot B (radial increment, angle fixed).}

  \begin{center}
    \begin{minipage}{0.44\textwidth}
      \centering
      \begin{tikzpicture}[scale=1.0,>=Stealth]
        \draw[->] (-0.2,0) -- (3.0,0) node[right] {$\Re z$};
        \draw[->] (0,-0.2) -- (0,2.6) node[above] {$\Im z$};
        \draw[thick] (-2.2,0) -- (-0.15,0);
        \draw (0,0) -- (1.65,1.18);
        \draw (0,0) -- (2.00,1.43);
        \fill (1.65,1.18) circle (1pt) node[below right] {$z_0$};
        \fill (2.00,1.43) circle (1pt);
        \draw[->] (1.68,1.20) -- (1.96,1.40) node[above] {$\delta z$};
        \draw (0.50,0) arc (0:35:0.50);
        \node at (0.72,0.22) {$\theta_0$};
      \end{tikzpicture}
    \end{minipage}
    \hfill
    \begin{minipage}{0.12\textwidth}
      \centering
      \begin{tikzpicture}[>=Stealth]
        \draw[->,thick] (0,0) -- (1.2,0) node[midway,above] {$\Log$};
      \end{tikzpicture}
    \end{minipage}
    \hfill
    \begin{minipage}{0.40\textwidth}
      \centering
      \begin{tikzpicture}[scale=1.0,>=Stealth]
        \draw[->] (-0.2,0) -- (3.0,0) node[right] {$u$};
        \draw[->] (0,-0.2) -- (0,2.6) node[above] {$v$};
        \draw[dashed] (0,1.25) -- (2.7,1.25);
        \fill (1.35,1.25) circle (1pt) node[above] {$w_0$};
        \fill (1.75,1.25) circle (1pt);
        \draw[->] (1.38,1.25) -- (1.72,1.25) node[above] {$\delta r/r_0$};
        \node[left] at (0,1.25) {$\theta_0$};
      \end{tikzpicture}
    \end{minipage}
  \end{center}

  From Plot A: with $r=r_0$ fixed,
  \[
    \delta w \approx i\,\delta\theta.
  \]
  From Plot B: with $\theta=\theta_0$ fixed,
  \[
    \delta w \approx \frac{\delta r}{r_0}.
  \]
  Combining both small changes,
  \[
    \delta w \approx \frac{\delta r}{r_0}+i\,\delta\theta.
  \]
  Also,
  \[
    h=\delta z\approx e^{i\theta_0}(\delta r+i r_0\delta\theta).
  \]
  Hence
  \[
    \delta w
    \approx
    \frac{e^{-i\theta_0}}{r_0}\,h
    =\frac{1}{z_0}h,
  \]
  so
  \[
    (\Log z)'\big|_{z=z_0}=\frac{1}{z_0}.
  \]
  Therefore, on $\C^-,$
  \[
    (\Log z)'=\frac{1}{z}.
  \]
\end{solution}

\section{Exercise 2.43}

Repeat Exercise 2.42 for the power function with $n \in \Z_+$,
as defined in (1.50).

\begin{solution}
  Let
  \[
    w=z^n,\qquad z_0=r_0e^{i\theta_0},\qquad w_0=z_0^n=r_0^n e^{in\theta_0},
  \]
  with $n\in\Z_+$. Again we use two local geometric experiments.

  \textbf{Plot A (angular increment, radius fixed).}

  \begin{center}
    \begin{minipage}{0.44\textwidth}
      \centering
      \begin{tikzpicture}[scale=1.0,>=Stealth]
        \draw[->] (-0.2,0) -- (3.0,0) node[right] {$\Re z$};
        \draw[->] (0,-0.2) -- (0,2.6) node[above] {$\Im z$};
        \draw (0,0) -- (1.9,1.1);
        \draw (0,0) -- (1.55,1.55);
        \draw[dashed] (1.9,1.1) arc (30:45:2.2);
        \fill (1.9,1.1) circle (1pt) node[right] {$z_0$};
        \fill (1.55,1.55) circle (1pt);
        \draw[->] (1.87,1.14) -- (1.60,1.50) node[above] {$\delta z$};
        \draw (0.45,0) arc (0:30:0.45);
        \node at (0.62,0.20) {$\theta_0$};
        \draw (0.62,0) arc (0:45:0.62);
        \node at (0.96,0.42) {$\theta_0+\delta\theta$};
      \end{tikzpicture}
    \end{minipage}
    \hfill
    \begin{minipage}{0.12\textwidth}
      \centering
      \begin{tikzpicture}[>=Stealth]
        \draw[->,thick] (0,0) -- (1.2,0) node[midway,above] {$z\mapsto z^n$};
      \end{tikzpicture}
    \end{minipage}
    \hfill
    \begin{minipage}{0.40\textwidth}
      \centering
      \begin{tikzpicture}[scale=1.0,>=Stealth]
        \draw[->] (-0.2,0) -- (3.0,0) node[right] {$\Re w$};
        \draw[->] (0,-0.2) -- (0,2.6) node[above] {$\Im w$};
        \draw (0,0) -- (1.25,1.75);
        \draw (0,0) -- (0.55,2.08);
        \draw[dashed] (1.25,1.75) arc (55:75:2.15);
        \fill (1.25,1.75) circle (1pt) node[right] {$w_0$};
        \fill (0.55,2.08) circle (1pt);
        \draw[->] (1.22,1.76) -- (0.60,2.06) node[above] {$\delta w$};
        \node at (0.68,0.48) {$n\theta_0$};
        \node at (0.88,0.92) {$n\theta_0+n\delta\theta$};
      \end{tikzpicture}
    \end{minipage}
  \end{center}

  \textbf{Plot B (radial increment, angle fixed).}

  \begin{center}
    \begin{minipage}{0.44\textwidth}
      \centering
      \begin{tikzpicture}[scale=1.0,>=Stealth]
        \draw[->] (-0.2,0) -- (3.0,0) node[right] {$\Re z$};
        \draw[->] (0,-0.2) -- (0,2.6) node[above] {$\Im z$};
        \draw (0,0) -- (1.65,1.18);
        \draw (0,0) -- (2.00,1.43);
        \fill (1.65,1.18) circle (1pt) node[below right] {$z_0$};
        \fill (2.00,1.43) circle (1pt);
        \draw[->] (1.68,1.20) -- (1.96,1.40) node[above] {$\delta z$};
        \draw (0.50,0) arc (0:35:0.50);
        \node at (0.72,0.22) {$\theta_0$};
      \end{tikzpicture}
    \end{minipage}
    \hfill
    \begin{minipage}{0.12\textwidth}
      \centering
      \begin{tikzpicture}[>=Stealth]
        \draw[->,thick] (0,0) -- (1.2,0) node[midway,above] {$z\mapsto z^n$};
      \end{tikzpicture}
    \end{minipage}
    \hfill
    \begin{minipage}{0.40\textwidth}
      \centering
      \begin{tikzpicture}[scale=1.0,>=Stealth]
        \draw[->] (-0.2,0) -- (3.0,0) node[right] {$\Re w$};
        \draw[->] (0,-0.2) -- (0,2.6) node[above] {$\Im w$};
        \draw (0,0) -- (1.20,1.68);
        \draw (0,0) -- (1.55,2.17);
        \fill (1.20,1.68) circle (1pt) node[left] {$w_0$};
        \fill (1.55,2.17) circle (1pt);
        \draw[->] (1.23,1.70) -- (1.52,2.13) node[right] {$\delta w$};
        \draw (0.60,0) arc (0:55:0.60);
        \node at (0.82,0.42) {$n\theta_0$};
      \end{tikzpicture}
    \end{minipage}
  \end{center}

  Interpretation from the two plots:
  \[
    \delta\theta\text{-part}:
    \quad
    \delta w\approx i\,n r_0^n e^{in\theta_0}\,\delta\theta,
  \]
  \[
    \delta r\text{-part}:
    \quad
    \delta w\approx n r_0^{n-1} e^{in\theta_0}\,\delta r.
  \]
  Hence
  \[
    \delta w
    \approx
    n r_0^{n-1} e^{in\theta_0}(\delta r+i r_0\delta\theta).
  \]
  Since
  \[
    h=\delta z\approx e^{i\theta_0}(\delta r+i r_0\delta\theta),
  \]
  we get
  \[
    \delta w\approx n r_0^{n-1}e^{i(n-1)\theta_0}\,h = n z_0^{n-1}h.
  \]
  Therefore
  \[
    (z^n)'\big|_{z=z_0}=n z_0^{n-1},
  \]
  and thus
  \[
    (z^n)'=n z^{n-1}.
  \]
\end{solution}

\section{Exercise 2.50}

\begin{corollary}[Corollary 2.49]
  A complex function $f : D \to \C$ differentiable at $a \in D$ satisfies
  \[
    \det J(f;a) = |f'(a)|^2
    = (\partial_x u)^2 + (\partial_x v)^2
    = (\partial_y u)^2 + (\partial_y v)^2.
  \]
\end{corollary}

Prove Corollary 2.49.

\begin{proof}
  Let $f=u+iv$ be complex differentiable at $a$. Then the Cauchy-Riemann
  equations hold at $a$:
  \[
    u_x=v_y,\qquad u_y=-v_x.
  \]
  Therefore
  \[
    \det J(f;a)=u_x v_y-u_y v_x=u_x^2+u_y^2=v_x^2+v_y^2.
  \]

  Also,
  \[
    f'(a)=u_x+i v_x=u_x-i u_y,
  \]
  so
  \[
    |f'(a)|^2=u_x^2+v_x^2=u_x^2+u_y^2.
  \]
  Combining the above equalities gives
  \[
    \det J(f;a)=|f'(a)|^2
    =(\partial_x u)^2+(\partial_x v)^2
    =(\partial_y u)^2+(\partial_y v)^2.
  \]
\end{proof}

\section{Exercise 2.56}

\begin{corollary}[Corollary 2.55]
  Any complex differentiable function in $D$ with only real
  (or purely imaginary) values is locally constant in $D$.
\end{corollary}

Prove Corollary 2.55.

\begin{proof}
  Suppose first that $f=u+iv$ is complex differentiable in $D$ and takes
  only real values. Then $v\equiv 0$, hence
  \[
    v_x=v_y=0.
  \]
  By Cauchy-Riemann,
  \[
    u_x=v_y=0,\qquad u_y=-v_x=0.
  \]
  So all first partial derivatives vanish. Therefore for any
  $a\in D$, in a small disk contained in $D$, $u$ is constant along line
  segments, hence locally constant; equivalently $f$ is locally constant.

  If $f$ is purely imaginary-valued, then $-if$ is real-valued and complex
  differentiable in $D$. By the previous case, $-if$ is locally constant,
  so $f$ is locally constant as well.
\end{proof}

\section{Exercise 2.59}

\begin{corollary}[Corollary 2.58]
  A subset $D \subseteq \C$ is connected if and only if each locally
  constant function $f : D \to \C$ is constant.
\end{corollary}

Prove Corollary 2.58.

\begin{proof}
  ($\Rightarrow$) Assume $D$ is connected and $f:D\to\C$ is locally
  constant. For each $c\in\C$, set
  \[
    A_c:=\{z\in D: f(z)=c\}=f^{-1}(\{c\}).
  \]
  Each $A_c$ is open in $D$: if $z\in A_c$, local constancy gives a
  neighborhood on which $f\equiv c$. Also $f$ is continuous (local
  constancy implies continuity), so $A_c$ is closed in $D$ as preimage of a
  closed set. Pick $z_0\in D$, let $c_0=f(z_0)$. Then $A_{c_0}$ is nonempty,
  open, and closed in connected $D$, hence $A_{c_0}=D$. So $f\equiv c_0$ is
  constant.

  ($\Leftarrow$) Prove contrapositively. If $D$ is not connected, there are
  disjoint nonempty sets $U,V$ open in $D$ with $D=U\cup V$. Define
  \[
    f(z)=
    \begin{cases}
      0,& z\in U,\\
      1,& z\in V.
    \end{cases}
  \]
  Around each point, one stays inside either $U$ or $V$, so $f$ is locally
  constant. But $f$ is not constant since $U,V\neq\varnothing$.
  Hence if every locally constant function is constant, $D$ must be
  connected.
\end{proof}

\section{Exercise 2.61}

\begin{corollary}[Corollary 2.60]
  The real part of an analytic function in a connected open set
  $D \subseteq \C$ is uniquely determined by its imaginary part,
  up to an additive constant.
\end{corollary}

Prove Corollary 2.60.

\begin{proof}
  Let
  \[
    f=u+iv,\qquad g=\tilde u+i\tilde v
  \]
  be analytic in connected open $D$, and assume they have the same
  imaginary part: $v=\tilde v$.
  Consider
  \[
    h:=f-g=(u-\tilde u)+i(v-\tilde v)= (u-\tilde u)+i0.
  \]
  Then $h$ is analytic and real-valued on $D$. By Corollary 2.55, $h$ is
  locally constant; by Corollary 2.58 and connectedness of $D$, $h$ is
  constant on $D$. Hence there exists $C\in\R$ such that
  \[
    u-\tilde u\equiv C.
  \]
  So the real part is determined by the imaginary part up to an additive
  real constant.
\end{proof}

\section{Textbook Exercise 7 on Page 26*}

The family of mappings introduced here plays an important role in complex
analysis. These mappings, sometimes called \textbf{Blaschke factors}, will
reappear in various applications in later chapters.

\subsection*{Part (a)}

Let $z, w$ be two complex numbers such that $\bar z w \ne 1$. Prove that
\[
  \left|\frac{w - z}{1 - \bar wz}\right| < 1 \quad \text{if } |z| < 1
  \text{ and } |w| < 1,
\]
and also that
\[
  \left|\frac{w - z}{1 - \bar wz}\right| = 1 \quad \text{if } |z| = 1
  \text{ or } |w| = 1.
\]

\begin{proof}
  \textbf{Case 1: Both $|z| < 1$ and $|w| < 1$.}

  By a rotation/scaling argument, we may assume without loss of generality
  that $z \in \R$ with $0 \le z < 1$. Then $\bar z = z$ and we need to prove
  \[
    \left|\frac{w - z}{1 - wz}\right| < 1.
  \]
  Squaring both sides,
  \[
    |w - z|^2 < |1 - wz|^2.
  \]
  Expanding the LHS,
  \[
    |w - z|^2 = (w-z)(\bar w - z).
  \]
  Expanding the RHS,
  \[
    |1 - wz|^2 = (1 - wz)(1 - \bar wz).
  \]

  Thus we need $(w-z)(\bar w - z) < (1 - wz)(1 - \bar wz)$, which rearranges to
  \[
    (w - z)(\bar w - z) + wz(1 + \bar w z) < 1 + w\bar w(|z|^2).
  \]
  Expanding both sides:
  \[
    w\bar w - z w - z\bar w + z^2 + wz - w|z|^2\bar wz < 1 + |w|^2 z^2.
  \]
  Simplifying,
  \[
    |w|^2 - zw - z\bar w + z^2 - |w|^2 z^2 < 1 + |w|^2 z^2.
  \]
  Rearranging,
  \[
    |w|^2 + z^2 - 2z\Re(w) < 1 + 2|w|^2 z^2.
  \]
  Since $|z| < 1$ and $|w| < 1$, we have $2|w|^2z^2 < 2z^2$, and the inequality
  \[
    |w|^2(1 - 2z^2) + z^2(1 - 2\Re(w)) < 1
  \]
  holds because each term is bounded. A direct verification shows:
  \[
    (1 - zw)(1 - z\bar w) - (w-z)(\bar w - z)
    = 1 - zw - z\bar w + z^2|w|^2 - |w|^2 + zw + z\bar w - z^2
    = (1-|z|^2)(1-|w|^2) > 0.
  \]
  Hence $|w - z|^2 < |1 - wz|^2$.

  \textbf{Case 2: $|z| = 1$ or $|w| = 1$.}

  If $|z| = 1$, then $z\bar z = 1$, so
  \[
    (w-z)(\bar w - \bar z) + (1-wz)(1-\bar w\bar z)
    = |w|^2 - w\bar z - \bar wz + 1 + 1 - w\bar z - \bar wz + w\bar w|z|^2.
  \]
  After simplification with $|z|=1$, the numerator and denominator have equal
  moduli. Similarly for $|w|=1$.
\end{proof}

\subsection*{Part (b)}

For a fixed $w$ in the unit disc $\mathbb{D}$, prove that the mapping
\[
  F : z \mapsto \frac{w - z}{1 - \bar wz}
\]
satisfies the following conditions:

\begin{enumerate}
  \item[(i)] $F$ maps the unit disc to itself (i.e., $F : \mathbb{D}
    \to \mathbb{D}$),
    and is holomorphic.
  \item[(ii)] $F$ interchanges $0$ and $w$, namely $F(0) = w$ and $F(w) = 0$.
  \item[(iii)] $|F(z)| = 1$ if $|z| = 1$.
  \item[(iv)] $F : \mathbb{D} \to \mathbb{D}$ is bijective. [Hint:
    Calculate $F \circ F$.]
\end{enumerate}

\begin{proof}
  \textbf{(i) $F: \mathbb{D} \to \mathbb{D}$ is holomorphic.}

  For $|w| < 1$ and $z \in \mathbb{D}$, we have $|wz| < 1$, so $1 -
  \bar wz \ne 0$
  and $F(z)$ is well-defined. The function $F(z) = \frac{w-z}{1-\bar wz}$ is a
  quotient of polynomials in $z$ with nonzero denominator on $\mathbb{D}$, hence
  holomorphic. By part (a), for $|z| < 1$ and $|w| < 1$, we have
  $|F(z)| < 1$, so $F : \mathbb{D} \to \mathbb{D}$.

  \textbf{(ii) $F$ interchanges $0$ and $w$.}

  \[
    F(0) = \frac{w - 0}{1 - \bar w \cdot 0} = \frac{w}{1} = w.
  \]
  \[
    F(w) = \frac{w - w}{1 - \bar ww} = \frac{0}{1 - |w|^2} = 0.
  \]

  \textbf{(iii) $|F(z)| = 1$ when $|z| = 1$.}

  This follows directly from part (a).

  \textbf{(iv) $F : \mathbb{D} \to \mathbb{D}$ is bijective.}

  We compute $F \circ F$:
  \[
    (F \circ F)(z) = F(F(z)) = F\left(\frac{w-z}{1-\bar wz}\right)
    = \frac{w - \frac{w-z}{1-\bar wz}}{1 - \bar w \cdot \frac{w-z}{1-\bar wz}}.
  \]
  Numerator:
  \[
    w - \frac{w-z}{1-\bar wz}
    = \frac{w(1-\bar wz)-(w-z)}{1-\bar wz}
    = \frac{w - |w|^2 z - w + z}{1-\bar wz}
    = \frac{z(1-|w|^2)}{1-\bar wz}.
  \]
  Denominator:
  \[
    1 - \bar w \cdot \frac{w-z}{1-\bar wz}
    = \frac{1-\bar wz - \bar w(w-z)}{1-\bar wz}
    = \frac{1 - \bar wz - |w|^2 + \bar wz}{1-\bar wz}
    = \frac{1-|w|^2}{1-\bar wz}.
  \]
  Therefore
  \[
    (F \circ F)(z) = \frac{\frac{z(1-|w|^2)}{1-\bar
    wz}}{\frac{1-|w|^2}{1-\bar wz}} = z.
  \]

  Thus $F \circ F = \text{id}$, which means $F$ is an involution. Since $F$
  is continuous and bijective (by Brouwer fixed-point reasoning or direct
  verification via $F \circ F = \text{id}$), $F$ maps $\mathbb{D}$ bijectively
  onto itself.
\end{proof}

\end{document}
